%%
%% This is file `sigconf.tex',
%% generated with the docstrip utility.
%%
%% The original source files were:
%%
%% samples.dtx  (with options: `all,proceedings,sigconf')
%% 
%% IMPORTANT NOTICE:
%% 
%% For the copyright see the source file.
%% 
%% Any modified versions of this file must be renamed
%% with new filenames distinct from sigconf.tex.
%% 
%% For distribution of the original source see the terms
%% for copying and modification in the file samples.dtx.
%% 
%% This generated file may be distributed as long as the
%% original source files, as listed above, are part of the
%% same distribution. (The sources need not necessarily be
%% in the same archive or directory.)
%%
%%
%% Commands for TeXCount
%TC:macro \cite [option:text,text]
%TC:macro \citep [option:text,text]
%TC:macro \citet [option:text,text]
%TC:envir table 0 1
%TC:envir table* 0 1
%TC:envir tabular [ignore] word
%TC:envir displaymath 0 word
%TC:envir math 0 word
%TC:envir comment 0 0
%%
%% The first command in your LaTeX source must be the \documentclass
%% command.
%%
%% For submission and review of your manuscript please change the
%% command to \documentclass[manuscript, screen, review]{acmart}.
%%
%% When submitting camera ready or to TAPS, please change the command
%% to \documentclass[sigconf]{acmart} or whichever template is required
%% for your publication.
%%
%%
\documentclass[sigconf]{acmart}
\setlength{\headheight}{15.5pt}
%%
%% \BibTeX command to typeset BibTeX logo in the docs
\AtBeginDocument{%
  \providecommand\BibTeX{{%
    Bib\TeX}}}

%% Rights management information.  This information is sent to you
%% when you complete the rights form.  These commands have SAMPLE
%% values in them; it is your responsibility as an author to replace
%% the commands and values with those provided to you when you
%% complete the rights form.
\setcopyright{none}
\copyrightyear{2026}
\acmYear{2026}
\settopmatter{printacmref=false}
\renewcommand\footnotetextcopyrightpermission[1]{}
%% These commands are for a PROCEEDINGS abstract or paper.
\acmConference[ICAIF '26]{7th ACM International Conference on AI in
  Finance}{November 14--17, 2026}{Milan, Italy}
%%
%% Submission ID.
%% Use this when submitting an article to a sponsored event. You'll
%% receive a unique submission ID from the organizers
%% of the event, and this ID should be used as the parameter to this command.
%%\acmSubmissionID{123-A56-BU3}

%%
%% For managing citations, it is recommended to use bibliography
%% files in BibTeX format.
%%
%% You can then either use BibTeX with the ACM-Reference-Format style,
%% or BibLaTeX with the acmnumeric or acmauthoryear sytles, that include
%% support for advanced citation of software artefact from the
%% biblatex-software package, also separately available on CTAN.
%%
%% Look at the sample-*-biblatex.tex files for templates showcasing
%% the biblatex styles.
%%
%%
%% The majority of ACM publications use numbered citations and
%% references.  The command \citestyle{authoryear} switches to the
%% "author year" style.
%%
%% If you are preparing content for an event
%% sponsored by ACM SIGGRAPH, you must use the "author year" style of
%% citations and references.
%% Uncommenting
%% the next command will enable that style.
%%\citestyle{acmauthoryear}


%%
%% end of the preamble, start of the body of the document source.
\begin{document}

%%
%% The "title" command has an optional parameter,
%% allowing the author to define a "short title" to be used in page headers.
\title[Empirical Evaluation of ProtoHedge]{Empirical Evaluation of
ProtoHedge:\\Interpretable Deep Hedging with Real Market Data}

%%
%% The "author" command and its associated commands are used to define
%% the authors and their affiliations.
%% Of note is the shared affiliation of the first two authors, and the
%% "authornote" and "authornotemark" commands
%% used to denote shared contribution to the research.

\author{Yanis Zaim}
\affiliation{%
  \institution{Columbia University}
  \city{New York}
  \state{NY}
  \country{USA}}
\email{ymz2108@columbia.edu}

\author{Sayak Chakrabarty}
\affiliation{%
  \institution{Northwestern University}
  \city{Evanston}
  \state{IL}
  \country{USA}}
\email{sayak.chakrabarty@northwestern.edu}

\author{Imon Banerjee}
\affiliation{%
  \institution{Northwestern University}
  \city{Evanston}
  \state{IL}
  \country{USA}}
\email{imon.banerjee@northwestern.edu}
%%
%% By default, the full list of authors will be used in the page
%% headers. Often, this list is too long, and will overlap
%% other information printed in the page headers. This command allows
%% the author to define a more concise list
%% of authors' names for this purpose.
\renewcommand{\shortauthors}{Zaim, Chakrabarty, and Banerjee}

%%
%% The abstract is a short summary of the work to be presented in the
%% article.
\begin{abstract}
Deep Hedging learns sequential hedging policies directly from market states,
transaction costs, and a risk-sensitive objective, but its neural decisions
are difficult to audit. ProtoHedge addresses this problem by routing each
hedge through a finite set of representative market states, or prototypes:
the action is determined by the current state's similarity to those
prototypes and their learned hedge responses. We test whether this
performance--interpretability tradeoff survives historical option data, where
frictions, non-stationarity, and path dependence matter. ProtoHedge is compared
with unhedged exposure, delta rules, and a common-protocol black-box Deep
Hedging benchmark on European and arithmetic-average Asian short calls. The
design uses chronological purged splits, one consistent listed hedge contract
option per episode, cumulative-inventory accounting, payoff-sufficient Asian
states, hard bounds, and validation selection. The outcome excludes the
liability premium and therefore is not complete economic profit and loss.
Across the current ten-underlying panel, validation-selected ProtoHedge
improves the average mean and lower-tail offset relative to Deep Hedging for
both liabilities. European policies also reduce average position-bound
occupancy; the Asian bound result is weaker. Prototype use is concentrated,
varies with market regimes, and can be traced to historical states. The
evidence supports ProtoHedge as an inspectable empirical hedge without claiming
uniform dominance across assets or objectives.
\end{abstract}

%%
%% The code below is generated by the tool at http://dl.acm.org/ccs.cfm.
%% Please copy and paste the code instead of the example below.
%%
\ccsdesc[500]{Applied computing~Finance}
\ccsdesc[300]{Computing methodologies~Instance-based learning}
\ccsdesc[300]{Computing methodologies~Neural networks}
\ccsdesc[300]{Computing methodologies~Sequential decision making}

\keywords{deep hedging, ProtoHedge, interpretable machine learning,
option hedging, empirical finance}



%%
%% This command processes the author and affiliation and title
%% information and builds the first part of the formatted document.
\maketitle


\section{Introduction}
Deep Hedging replaces a prescribed analytical hedge with a policy learned
from market states, transaction costs, and a risk-sensitive objective. This
flexibility is valuable when hedging is discrete, costly, constrained, or
exposed to dynamics that depart from a stylized pricing model
\cite{buehler2019deep}. Its central
practical weakness is auditability. Strong out-of-sample performance alone
does not reveal which market states drive a hedge, whether comparable states
receive comparable actions, or whether the policy has learned a stable
economic rule rather than an opaque statistical association.

ProtoHedge addresses this limitation by routing each hedge decision through a
finite set of representative market states, called prototypes. A current
state is compared with these prototypes, and the resulting similarity weights
combine prototype-associated hedge responses into the final action. This
construction makes each action traceable to observed representative states
while retaining the risk-sensitive learning objective of Deep Hedging. The
original ProtoHedge study established architectural feasibility in
Black--Scholes and stochastic-volatility simulations. It did not establish
whether the same performance--interpretability tradeoff survives historical
non-stationarity, changing liquidity and volatility, observed option paths,
and path-dependent liabilities \cite{faloughi2025protohedge}.

This paper addresses that empirical gap using historical U.S.\ equity and ETF
option data. We preserve the original fixed-prototype architecture and compare
it with transparent delta-style rules and a black-box Deep Hedging benchmark
under a common state space, action space, objective, cost model, and position
constraints. The evaluation uses chronological train, validation, and test
periods; constructs rolling episodes only after the calendar is split; follows
one listed hedge-option contract throughout each episode; selects every model
configuration without test access; and evaluates uncertainty with paired
block-bootstrap intervals. We study both a terminal European call and an
arithmetic-average Asian call whose policy state includes the running payoff
state.

Under the current empirical results, validation-selected ProtoHedge policies
remain competitive with the Deep Hedging benchmark across both liabilities.
On the European panel, separate mean- and tail-oriented selections improve the
corresponding average liability-offset criteria while contacting hard position
bounds less frequently on average. The Asian experiment preserves the mean and
lower-tail competitiveness of ProtoHedge, although its position-bound advantage
is weaker, showing that the tradeoff changes under path dependence. Prototype
usage is concentrated, changes across return, volatility, and drawdown regimes,
and can be traced to concrete historical states. These findings support a
qualified conclusion: the fixed prototype bottleneck can preserve empirical
hedging quality while making decisions more inspectable, but its advantage is
not uniform across liabilities, assets, or evaluation criteria.

The paper makes four contributions:
\begin{enumerate}
    \item We construct 35,987 contract-consistent 20-step episodes for ten
    U.S.\ equity and ETF underlyings using WRDS spot and option data.
    \item We evaluate European and Asian short-call liabilities under
    inventory-based accounting, stylized transaction costs, hard position
    bounds, and a common-protocol comparison with classical and neural
    benchmarks.
    \item We split by calendar, select configurations and checkpoints using
    validation data only, and apply paired circular-block inference.
    \item We test interpretability through exact prototype-weighted action
    structure, prototype concentration, regime-dependent usage, and retrieval
    of nearest historical episodes.
\end{enumerate}









\section{ProtoHedge and the Empirical Validation Problem}
\label{sec:protohedge_empirical_problem}

\subsection{Deep Hedging as Sequential Risk Management}
\label{subsec:deep_hedging_background}

Deep hedging replaces a prescribed replication formula with a policy learned
from observed market states and a risk-sensitive hedging objective. Let
$P_t \in \mathbb{R}^m$ be the vector of hedge-instrument prices, let $h_t$ be
the cumulative hedge inventory, and let $\tau_t$ denote time remaining. At
time $t$, a policy observes a state such as
\[
x_t = (P_t, h_{t-1}, \tau_t)
\]
and outputs a trade increment
\[
a_t = \pi_\theta(x_t), \qquad
h_t = \Pi_{\mathcal{H}}(h_{t-1}+a_t),
\]
where $\Pi_{\mathcal{H}}$ enforces the admissible position set. Let
$H(P_{0:T})\leq 0$ be the terminal cash flow of the short liability and let
$\kappa_t$ be the per-unit transaction-cost vector. The premium-excluding
liability offset can be written schematically as
\[
\mathcal{O}^{\pi_\theta}
=
H(P_{0:T})
+
\sum_{t=0}^{T-1} h_t^\top(P_{t+1}-P_t)
-
\sum_{t=0}^{T-1}\kappa_t^\top |a_t|.
\]
The agent is trained to maximize the optimized-certainty-equivalent monetary
utility used in the original Deep Hedging framework. For the piecewise-linear
utility selected in our experiments,
\[
\mathcal{U}_{\lambda}(\mathcal{O})
=
\sup_{y\in\mathbb{R}}
\mathbb{E}\!\left[
(1+\lambda)\min(\mathcal{O}+y,0)-y
\right],
\qquad \lambda=1.
\]
The scalar $y$ is learned jointly with the policy. We call this implementation
setting the \texttt{cvar} monetary utility because that is its name in the
original software; it must not be confused with the separately reported
empirical CVaR$_{5\%}$. The policy optimizes this risk-sensitive objective
rather than reproducing a model-implied delta exactly. This formulation
accommodates discrete rebalancing, costs, hard bounds, and dynamics that need
not follow a stylized pricing model \cite{buehler2019deep}. Because
$\mathcal{O}$ excludes the
initial premium received for selling the liability, it is not complete
economic profit and loss; its purpose is to compare how effectively policies
offset the same liability under a common objective.

A standard, or black-box, Deep Hedging policy represents $\pi_\theta$ by a
feed-forward neural network. It is a necessary performance benchmark in this
paper because it has the same information, objective, costs, and constraints
as ProtoHedge. Its limitation is diagnostic rather than computational: an
action can be evaluated without revealing which market states motivated it or
whether the policy has learned an economically stable rule.

\subsection{ProtoHedge: Fixed Prototypes as an Interpretable Policy}
\label{subsec:protohedge_background}

ProtoHedge retains the deep-hedging objective but restricts action generation
to a finite collection of representative states
\[
\mathcal{P}=\{p_1,\ldots,p_K\}.
\]
In the fixed-prototype design studied here, each $p_k$ is an observed
standardized state selected from the empirical state distribution. For a
standardized current state $\tilde{x}_t$, nonnegative feature weights
$\gamma$, and a shared learned projection $f_\phi$, ProtoHedge forms
\[
z_t=f_\phi(\gamma\odot\tilde{x}_t), \qquad
z_k^{p}=f_\phi(\gamma\odot p_k),
\]
then measures squared distances and converts them into normalized prototype
weights:
\[
d_{t,k}=\lVert z_t-z_k^{p}\rVert_2^2,
\qquad
\alpha_{t,k}
=
\frac{\exp(-d_{t,k})}
{\sum_{j=1}^{K}\exp(-d_{t,j})}.
\]
Each prototype has an associated trainable, bounded hedge response
$b_k\in\mathbb{R}^{m}$. The action is the similarity-weighted combination
\[
a_t=\sum_{k=1}^{K}\alpha_{t,k}b_k.
\]

This decomposition is the central ProtoHedge mechanism. For every hedge
decision, the model exposes the prototypes with the largest weights, their
representative market states, and their associated actions. A reviewer can
therefore trace a decision to examples of historically observed states rather
than only to opaque network parameters. Concentrated use of a small subset of
prototypes, systematic shifts in $\alpha_{t,k}$ across regimes, and retrieval
of the nearest historical episodes provide direct, testable forms of
interpretability \cite{faloughi2025protohedge}.

The anchors $p_k$ remain fixed during training in the present study; the
projection $f_\phi$ and the prototype actions $b_k$ are learned. This
preserves the original ProtoHedge premise that a prototype should remain tied
to an actual representative state, rather than drifting to an uninterpretable
parameter value. Section~\ref{sec:experimental_framework} specifies how the
prototypes are extracted, how $K$ and similarity weighting are varied, and how
the policy is trained.

\subsection{Why the Empirical Test Matters}
\label{subsec:empirical_gap}

The original ProtoHedge study showed that prototype-based policies can remain
competitive in Black--Scholes and stochastic volatility simulations
\cite{faloughi2025protohedge}. Historical
option data pose a materially different test: volatility and liquidity vary
over time, option observations are noisy, trading is costly and constrained,
and market paths contain regime shifts and discontinuities. The
interpretability claim is also stronger on real data, because a useful
prototype should correspond to a recognizable historical state rather than
only to a simulator-generated configuration.

Accordingly, this paper evaluates whether ProtoHedge remains competitive with
common-protocol Deep Hedging and transparent delta-style baselines while
avoiding excessive reliance on hard position bounds. It also tests whether
prototype use is concentrated, regime-dependent, and retrievable as concrete
historical episodes. We conduct this evaluation for both terminal European
short calls and path-dependent Asian short calls, allowing the
performance--interpretability tradeoff to be assessed under two economically
distinct liabilities.












\section{Experimental Framework}
\label{sec:experimental_framework}

This section describes the empirical design used to evaluate ProtoHedge on historical market data. Our implementation is deliberately conservative. We use fixed-length out-of-sample hedging episodes, explicit transaction costs, bounded trading actions, bounded cumulative hedge positions, validation-only model selection, and multi-seed evaluation. The goal is not merely to optimize a policy in sample, but to assess whether ProtoHedge remains competitive with Deep Hedging under a realistic and reproducible real-data protocol.

\subsection{Data and Episode Construction}
\label{subsec:data_episode_construction}

Our real-data experiments are based on historical U.S. equity and ETF option data obtained through WRDS, combining underlying spot-price histories with daily listed option records. The final panel contains ten liquid underlyings spanning broad equity indices, sector ETFs, fixed income, and single-name equities:
\[
\{\texttt{AAPL}, \texttt{IWM}, \texttt{NVDA}, \texttt{QQQ}, \texttt{SPY}, \texttt{TLT}, \texttt{XLE}, \texttt{XLF}, \texttt{XLK}, \texttt{XLV}\}.
\]
The available histories begin between January and February 2005, depending on
the ticker, and end in December 2024.

The raw OptionMetrics chain is restricted to American-style calls with finite
vendor Greeks and implied volatility, positive bids, non-crossed offers, and
sufficient initial liquidity, and is aligned with the corresponding CRSP
daily spot series. Midquotes define the listed option price. We split the
calendar before creating any
rolling window: training observations end on December 31, 2018, validation
observations run through December 31, 2021, and test observations begin in
2022. A 19-observation embargo, equal to the overlap of adjacent 20-step
windows, is removed after each boundary. Windows are then generated
independently inside each period, so no date can occur in more than one of the
training, validation, and test sets.

For every eligible 20-step spot window, we select the listed hedge option on
the first date according to:
\begin{enumerate}
    \item absolute distance between strike and spot,
    \item absolute distance between the contract's days-to-expiration and a 30-day target maturity,
    \item bid--ask spread,
    \item open interest (descending), and
    \item volume (descending),
\end{enumerate}
subject to an initial maturity filter of 20 to 40 calendar days, positive open
interest, relative bid--ask spread no greater than one, and absolute moneyness
no greater than 10\%. Crucially, we retain an episode only if the same
\texttt{optionid}, expiration, and strike remain valid on every episode date
and the contract expires after the liability horizon. Thus
\[
C_{i,t}=C(j_i,t), \qquad t=0,\ldots,T,
\]
for one fixed listed contract $j_i$ per episode. The option hedge is distinct
from the synthetic European or Asian liability, and no unmodeled option roll
occurs inside an episode.

The resulting per-date feature vector is
\[
x_t^{\mathrm{raw}} = \bigl(S_t,\; C_t,\; \Delta^{\mathrm{call}}_t,\; \nu_t,\; \sigma^{\mathrm{impl}}_t \bigr),
\]
where $S_t$ is the underlying spot, $C_t$ is the selected call price, $\Delta^{\mathrm{call}}_t$ is the selected call delta, $\nu_t$ is the selected call vega, and $\sigma^{\mathrm{impl}}_t$ is implied volatility. These features are stored in the episode tensor in the fixed order
\[
(\texttt{spot},\; \texttt{call\_price},\; \texttt{call\_delta},\; \texttt{call\_vega},\; \texttt{ivol}).
\]

Historical hedging episodes are constructed as rolling 20-step windows within
each chronological period. A new block begins whenever consecutive
observations are more than four calendar days apart, preventing windows from
joining disjoint histories. Letting $\tau=20$ denote the episode length, each
sample is a tensor in
\[
\mathbb{R}^{\tau \times 5}.
\]
Across the ten-ticker panel, the corrected contract-consistent procedure yields
35{,}987 episodes: 21{,}996 training, 6{,}955 validation, and 7{,}036 test
episodes.

\begin{table}[t]
\centering
\small
\begin{tabular}{lrrr}
\toprule
Ticker & Train & Validation & Test \\
\midrule
AAPL & 2,230 & 693 & 714 \\
IWM  & 2,501 & 701 & 715 \\
NVDA & 1,944 & 678 & 659 \\
QQQ  & 2,491 & 716 & 695 \\
SPY  & 2,531 & 718 & 715 \\
TLT  & 2,156 & 713 & 714 \\
XLE  & 2,369 & 655 & 713 \\
XLF  & 2,100 & 691 & 713 \\
XLK  & 1,781 & 704 & 707 \\
XLV  & 1,893 & 686 & 691 \\
\midrule
Total & 21,996 & 6,955 & 7,036 \\
\bottomrule
\end{tabular}
\caption{Contract-consistent 20-step episodes after the chronological split
and 19-observation boundary embargo. Windows may overlap within a period but
never across train, validation, and test periods.}
\label{tab:panel-counts}
\end{table}

Although the raw episode tensor contains five market features, the European
policies consume the state used by the Deep Hedging gym:
\[
(\texttt{price},\; \texttt{delta},\; \texttt{time\_left}),
\]
where \texttt{price} contains current spot and the price of the fixed listed
call, \texttt{delta} is the current cumulative inventory in both hedge
instruments, and \texttt{time\_left} is time to the liability horizon. For the
Asian liability, the policy additionally observes running-average moneyness,
\[
m^{\mathrm{Asian}}_{i,t}
=
\frac{\bar S_{i,t}}{K_i},
\qquad
\bar S_{i,t}=\frac{1}{t+1}\sum_{u=0}^{t}S_{i,u}.
\]
The same feature is included in Deep Hedging, ProtoHedge, prototype
clustering, similarity calculation, and prototype interpretation. This
prevents two paths with identical current prices and inventories but different
accumulated averages from being treated as the same liability state.

Finally, for numerical stability and cross-path comparability, the real-data world normalizes spot, call price, and call vega by each path's initial spot level. Under this normalization, a liability struck at-the-money at inception has normalized strike one, while remaining equivalent to a strike of $S_0$ in raw price units.

\subsection{Hedging Tasks}
\label{subsec:hedging_tasks}

We study two liabilities on the same historical episode panel.

\paragraph{European ATM short call.}
The primary task is to hedge a short at-the-money European call over a 20-step historical episode. If $S_{i,0}$ denotes the initial spot on episode $i$, then the strike is set to
\[
K_i = S_{i,0}.
\]
The liability payoff at maturity is
\[
H_i^{\mathrm{Euro}} = -\max(S_{i,T} - K_i, 0),
\]
where the negative sign reflects the fact that the hedger is short the option.

\paragraph{Asian ATM short call.}
The path-dependent extension replaces the European liability with an arithmetic-average Asian call, again struck at-the-money at episode start:
\[
K_i = S_{i,0}.
\]
Let the episode contain $T$ trading decisions. The Asian liability is
\[
H_i^{\mathrm{Asian}} = -\max\!\left(\frac{1}{T}\sum_{t=1}^{T} S_{i,t} - K_i,\, 0\right).
\]
In the implementation used for the final paper sweep, the averaging window spans the full historical episode.

\paragraph{Hedge instruments and actions.}
The hedge may trade in two liquid instruments:
\begin{enumerate}
    \item the underlying spot instrument, and
    \item the selected near-ATM listed call option.
\end{enumerate}
The action at each step is therefore two-dimensional,
\[
a_t = \bigl(a_t^{(S)}, a_t^{(C)}\bigr),
\]
where each component denotes a trade increment rather than an absolute target position. The cumulative hedge inventory evolves recursively as
\[
\delta_t = \delta_{t-1} + a_t.
\]

\paragraph{Step-wise hedge accounting.}
All reported experiments use step-wise inventory accounting. If
$\delta_{i,t}^{(S)}$ and $\delta_{i,t}^{(C)}$ are the cumulative positions
held after trading at $t$, and $\Delta S_{i,t}=S_{i,t+1}-S_{i,t}$ and
$\Delta C_{i,t}=C(j_i,t+1)-C(j_i,t)$ are one-step changes, then the trading
gain is
\[
\mathrm{TG}_i
=
\sum_{t=0}^{T-1}
\left(
\delta_{i,t}^{(S)}\Delta S_{i,t}
+
\delta_{i,t}^{(C)}\Delta C_{i,t}
\right).
\]
The use of cumulative inventory is essential because the real-data world
provides one-step price changes. It is algebraically equivalent to the
return-to-terminal accounting in the original synthetic ProtoHedge world, but
it does not incorrectly discard positions established at earlier dates.

\paragraph{Transaction costs.}
We impose the stylized transaction-cost specification inherited from the
experimental framework separately on spot and option trades. It is a
deterministic friction proxy rather than reconstructed executable bid--ask
slippage. The per-unit spot trading-cost coefficient is
\[
c_{i,t}^{(S)} = 0.0002\,|S_{i,t}|.
\]
The per-unit option trading-cost coefficient is
\[
c_{i,t}^{(C)} = 0.02\,|\nu_{i,t}| + 0.0002\,|\Delta^{\mathrm{call}}_{i,t}| + 0.0005\,|C_{i,t}|.
\]
The realized trading cost on episode $i$ is therefore
\[
\mathrm{Cost}_i
=
\sum_{t=0}^{T-1}
\left(
|a_{i,t}^{(S)}|\,c_{i,t}^{(S)}
+
|a_{i,t}^{(C)}|\,c_{i,t}^{(C)}
\right).
\]

\paragraph{Liability offset and premium convention.}
For both liabilities, the empirical outcome is
\[
\mathcal{O}_i = H_i + \mathrm{TG}_i - \mathrm{Cost}_i.
\]
The initial premium of the synthetic liability is not observed. In
particular, the listed call used as a hedge is a distinct contract and its
initial quote is not the premium of the European or Asian liability. We
therefore do not add that quote to $\mathcal{O}_i$ and do not call
$\mathcal{O}_i$ economic profit and loss. Higher values instead mean that the
trading strategy offsets more of the common short liability after costs. A
complete terminal P\&L would require an independently specified liability
premium $V_{i,0}$:
\[
\mathrm{P\&L}^{\mathrm{terminal}}_i
=
V_{i,0}+\mathcal{O}_i.
\]
All model comparisons in this paper use $\mathcal{O}_i$ consistently.

\subsection{Models and Baselines}
\label{subsec:models_baselines}

We compare ProtoHedge against both classical baselines and a black-box neural benchmark.

\paragraph{Unhedged.}
This baseline takes no hedge action:
\[
a_t \equiv 0.
\]
Its liability offset is simply the liability payoff.

\paragraph{Spot-Delta.}
This is a rule-based baseline that hedges only with the underlying spot instrument. At each step, the target cumulative spot position is set equal to the observed delta of the selected near-ATM call. The option hedge position remains zero. This baseline uses the market-provided option delta directly and therefore provides a transparent classical reference point.

\paragraph{Spot-Delta Band.}
This baseline extends Spot-Delta with a no-trade band. Instead of rebalancing to the point target at every step, the current cumulative spot position is left unchanged as long as it remains within
\[
[\Delta^{\mathrm{call}}_t - b,\; \Delta^{\mathrm{call}}_t + b],
\]
where $b \ge 0$ is a band width. If the current position exits the band, it is projected back to the closest boundary. We tune $b$ on the validation set using the grid
\[
b \in \{0.00,\;0.01,\;0.02,\;0.05,\;0.10,\;0.15\},
\]
selecting the band that maximizes validation mean liability offset.
This provides a friction-aware transparent comparator in the spirit of
classical no-transaction-region hedging \cite{whalley1997}.

\paragraph{Deep Hedging benchmark.}
The main black-box benchmark is a feed-forward Deep Hedging policy trained end-to-end under the same world, friction model, and objective as ProtoHedge. For the European task, the agent uses
\[
(\texttt{price},\; \texttt{delta},\; \texttt{time\_left}),
\]
and outputs a two-dimensional trade increment. The Asian version adds
$m_t^{\mathrm{Asian}}$. In the panel runs, the benchmark network has width 20,
depth 3, and Softplus activations.

\paragraph{ProtoHedge.}
ProtoHedge uses the same task-specific state representation and output space as
the Deep Hedging benchmark, but constrains action generation through a
prototype layer. In the full grid, we vary:
\begin{enumerate}
    \item the number of prototypes
    \[
    K \in \{10, 25, 50, 100\},
    \]
    \item the prototype source policy, using either the trained Deep Hedging
    benchmark or the Spot-Delta baseline as the source from which prototypes
    are extracted, and
    \item whether similarity is weighted or unweighted.
\end{enumerate}

Prototype extraction follows the same reusable pipeline throughout the repo. For a given source policy, we:
\begin{enumerate}
    \item construct the flattened feature matrix seen by the policy,
    \item standardize it with a \texttt{StandardScaler},
    \item apply $K$-means clustering in standardized feature space
    \cite{macqueen1967}, and
    \item store, for each cluster, the nearest real standardized feature vector rather than the raw cluster centroid.
\end{enumerate}
This last step matters because it keeps each stored prototype tied to an actually observed market state rather than to an abstract average point with no direct empirical counterpart.

The shared projection $f_\phi$ is a two-layer map
\[
D\longrightarrow 32\longrightarrow D
\]
with a ReLU between layers. Similarity uses unit temperature. In the
unweighted specification, all feature-distance weights equal one. In the
weighted specification, time remaining receives weight five and all other
state components receive weight one. These weights and the prototypes remain
fixed; the projection and bounded prototype actions are trainable. The action
bounds are implemented with the same smooth soft-clipping logic as the
original ProtoHedge layer. This keeps the present paper focused on empirical
validation of the original fixed-anchor design rather than the separate
question of moving prototypes.

\subsection{Training and Model Selection Protocol}
\label{subsec:training_protocol}

All real-data neural models are trained under the same multi-seed protocol.

\paragraph{Train/validation/test splits.}
The split is fixed by calendar and is identical across random seeds:
\[
\begin{aligned}
\text{training:} &\quad \text{through 2018-12-31},\\
\text{validation:} &\quad \text{2019 through 2021},\\
\text{test:} &\quad \text{2022 through 2024}.
\end{aligned}
\]
The calendar is split before windows are constructed, and the 19-observation
embargo at each boundary prevents a 20-step window from crossing periods.
Random seeds therefore affect neural initialization and prototype clustering,
not test membership. The main panel uses the three seeds
\[
(1234,\;2345,\;3456).
\]

\paragraph{Optimization.}
All neural models are trained for 800 epochs using the same
optimized-certainty-equivalent \texttt{cvar} monetary utility described in
Section~\ref{subsec:deep_hedging_background}, with risk-aversion parameter
one. This training objective is distinct from the empirical 5\% tail metric
used for evaluation. The optimizer uses learning rate
\[
10^{-3},
\]
together with clipping safeguards already built into the training pipeline. Classical baselines are not trained in the neural sense, but they are evaluated on exactly the same splits and under the same friction model.

\paragraph{Feature normalization and state semantics.}
The real-data world is run in normalized mode. Spot, call price, and call vega are divided pathwise by the initial spot. European policies observe current hedgeable prices, current hedge inventory, and time remaining; Asian policies also observe running-average moneyness. This keeps the neural state comparable across underlyings while supplying the payoff-relevant state of each liability.

\paragraph{Validation-based checkpoint selection.}
Model selection is based on the validation split, not on the training loss alone. Let $L_{\mathrm{val}}$ denote the validation objective. For the final robust panel runs, checkpoints are ranked by the selection score
\[
\mathcal{S}
=
L_{\mathrm{val}}
+
\alpha_a \,\overline{|a|}
+
\alpha_\delta \,\overline{|\delta|}
+
\alpha_b \,\mathrm{BO}
+
\alpha_p \,\mathrm{PT},
\]
where $\overline{|a|}$ is mean absolute action magnitude on the validation set, $\overline{|\delta|}$ is mean absolute cumulative position magnitude, $\mathrm{BO}$ is validation bound occupancy, and $\mathrm{PT}$ is the fraction of validation paths that touch any cumulative bound. In the reported full runs, we use
\[
\alpha_a = 0.005,\qquad
\alpha_\delta = 0.010,\qquad
\alpha_b = 0.100,\qquad
\alpha_p = 0.
\]
Thus, validation selection explicitly favors models that do not rely excessively on large actions, large cumulative positions, or persistent boundary saturation.

\paragraph{Validation-only configuration selection.}
Checkpoint selection and architecture selection are separate. First, every
candidate is trained on the training period, and its checkpoint is chosen by
the validation score above. Second, candidate-level validation metrics are
averaged across the three seeds. We apply the prespecified validation bound
screen and freeze two prespecified selections:
\begin{enumerate}
    \item \emph{Validation-Mean}: maximize validation mean liability offset,
    with validation CVaR, lower bound occupancy, smaller $K$, and model name as
    deterministic tie-breakers;
    \item \emph{Validation-Tail}: maximize validation
    CVaR$_{5\%}$ of liability offset, with validation mean, lower bound
    occupancy, smaller $K$, and model name as tie-breakers.
\end{enumerate}
This selection covers $K$, weighted versus unweighted similarity, and
prototype-source policy. Regularization coefficients, action bounds, position
bounds, selection objectives, and the validation robustness threshold are
prespecified before the test stage. Candidate ProtoHedge models are not
evaluated on test. After both selectors are written to disk, the baselines,
the validation-selected Deep Hedging checkpoint, and only the frozen
ProtoHedge configurations are evaluated once on the test period.

\paragraph{Regularization.}
In addition to validation-time selection penalties, the training loss includes mild regularization terms on action and cumulative-position magnitude:
\[
\lambda_a = 0.001,
\qquad
\lambda_\delta = 0.002.
\]
Concretely, if $L_{\mathrm{hedge}}$ is the base objective, then the optimized loss is
\[
L_{\mathrm{train}}
=
L_{\mathrm{hedge}}
+
\lambda_a \,\mathbb{E}[|a|]
+
\lambda_\delta \,\mathbb{E}[|\delta|].
\]
These penalties are deliberately mild: they do not dominate the objective, but they discourage the model from achieving performance through unnecessarily aggressive trading.

\paragraph{Trade bounds and cumulative position bounds.}
The final reported real-data runs impose both action-level and cumulative-position bounds. At every step and for each traded instrument, trade increments are clipped to
\[
[-1,\,1],
\]
and cumulative hedge positions are also constrained to
\[
[-1,\,1].
\]
These bounds are scientifically important. In exploratory unconstrained runs, apparent performance improvements could be accompanied by heavy saturation at the position boundaries. The bounded protocol is designed to test whether ProtoHedge remains competitive without relying on unrealistic edge-of-feasible-region behavior.

\paragraph{Paper sweep design.}
For each ticker and liability, the full sweep evaluates the complete
ProtoHedge grid across all three seeds on training and validation data.
Selection is performed separately by ticker and liability using the fixed
rules above. Only the resulting validation-selected configurations enter that
ticker's test comparison.

\begin{table}[t]
\centering
\small
\begin{tabular}{ll}
\toprule
Component & Prespecified value \\
\midrule
Episode length & 20 observations \\
Train / validation / test & 2005--18 / 2019--21 / 2022--24 \\
Boundary embargo & 19 observations \\
Random seeds & 1234, 2345, 3456 \\
Neural epochs & 800 \\
Learning rate & $10^{-3}$ \\
Deep Hedging network & width 20, depth 3, Softplus \\
Prototype counts & 10, 25, 50, 100 \\
Prototype sources & Deep Hedging, Spot-Delta \\
Similarity & weighted, unweighted \\
Trade / position bounds & $[-1,1]$ per instrument \\
Validation occupancy screen & $\mathrm{BO}\leq 0.50$ \\
\bottomrule
\end{tabular}
\caption{Common experimental protocol. Hyperparameter and checkpoint choices
are made on training and validation data; the test calendar is fixed across
seeds.}
\label{tab:protocol}
\end{table}

\subsection{Evaluation Metrics}
\label{subsec:evaluation_metrics}

We evaluate models at both the hedging-performance level and the interpretability level.

\paragraph{Primary performance metrics.}
For each fitted model and each test split, we compute:
\begin{enumerate}
    \item mean liability offset,
    \[
    \mathbb{E}[\mathcal{O}],
    \]
    \item empirical lower-tail Value-at-Risk and Conditional Value-at-Risk,
    \[
    \mathrm{VaR}_{5\%}(\mathcal{O}), \qquad
    \mathrm{CVaR}_{5\%}(\mathcal{O}),
    \]
    \item root mean squared liability offset,
    \[
    \mathrm{RMSE}(\mathcal{O})
    =
    \sqrt{\mathbb{E}[\mathcal{O}^2]},
    \]
    \item downside deviation,
    \[
    \mathrm{DD}(\mathcal{O})
    =
    \sqrt{\mathbb{E}[\min(\mathcal{O},0)^2]},
    \]
    and realized mean trading-cost and trading-gain summaries.
\end{enumerate}
Higher mean and CVaR are better; lower RMSE and downside deviation are better.
We use the lower-tail empirical expected-shortfall convention for gains
\cite{rockafellar2000}.
We record the negative-offset rate only as an implementation diagnostic. We
do not present it as shortfall probability or a headline economic metric:
because the premium is excluded, zero is not a break-even P\&L threshold, and
an unhedged out-of-the-money liability mechanically produces many zeros while
a hedge incurs trading costs.

\paragraph{Position-bound diagnostics.}
Because constrained hedging behavior is a central part of the empirical question, we also record:
\begin{enumerate}
    \item the fraction of all path--step--instrument positions lying at either cumulative position bound,
    \item the fraction lying specifically at the upper or lower bound, and
    \item the fraction of paths that touch any cumulative position bound at least once.
\end{enumerate}
We refer to the first quantity as \emph{bound occupancy}. Lower bound occupancy is interpreted as evidence that a model is operating more consistently in the interior of the feasible region, which is economically preferable all else equal.

\paragraph{Robust-screen criterion.}
The prespecified screen is calculated on validation data only. A candidate
passes when average validation bound occupancy is no greater than 0.50; the
path-touch threshold is recorded but not binding in the reported setup. Test
bound occupancy remains a diagnostic outcome and cannot determine whether a
model enters the paper comparison.

\paragraph{Dependence-aware uncertainty.}
Adjacent episodes within a period may share 19 of 20 dates, so individual
windows are not treated as independent observations. For each
validation-selected ProtoHedge policy and its matched Deep Hedging benchmark,
we first average the pathwise test offsets across seeds and then construct
paired 95\% circular-block bootstrap intervals with block length 20 and 2,000
replicates. We report intervals for differences in mean offset and
CVaR$_{5\%}$ and for improvements in RMSE and downside deviation; all are
oriented so positive values favor ProtoHedge. Panel intervals independently
block-resample within ticker, calculate each ticker statistic, and average
with equal ticker weight \cite{politis1992}.

\paragraph{Prototype concentration metrics.}
For ProtoHedge models, we further compute:
\begin{enumerate}
    \item average prototype weight across all paths and time steps,
    \item top-1 prototype frequency, i.e.\ the proportion of decisions for which each prototype receives the highest weight, and
    \item cumulative usage mass of the top-$k$ prototypes.
\end{enumerate}
These metrics quantify whether decisions are routed through a small number of dominant prototype states or diffusely across many prototypes.

\paragraph{Regime-sensitive interpretability metrics.}
To test whether prototype use changes systematically with market conditions, we form episode-level regime labels from the historical paths:
\begin{itemize}
    \item \emph{return regime}, using terciles of terminal return,
    \item \emph{volatility regime}, using terciles of average implied volatility, and
    \item \emph{drawdown regime}, using terciles of maximum realized drawdown.
\end{itemize}
This yields labels of the form
\[
\begin{aligned}
&\{\text{down},\text{middle},\text{up}\},\\
&\{\text{low\_vol},\text{mid\_vol},\text{high\_vol}\},\\
&\{\text{mild\_dd},\text{mid\_dd},\text{deep\_dd}\}.
\end{aligned}
\]
Within each regime, we summarize mean offset, tail risk, RMSE, downside
deviation, and average prototype weights. This lets us test whether the same
few prototypes dominate everywhere, or whether ProtoHedge reallocates
attention across market states in a regime-sensitive way.

\paragraph{Nearest historical states.}
Finally, for the dominant prototypes, we recover the nearest historical feature vectors and corresponding underlying paths. This allows us to move beyond abstract prototype coordinates and ask whether the most-used prototypes correspond to recognizable empirical market scenarios. This step is central to the interpretability claim of the paper: the prototype mechanism is not merely sparse, but empirically inspectable.










\section{Real-Data Empirical Results}
\label{sec:empirical_results}

We now examine whether the fixed-prototype restriction preserves hedging
quality on historical data and whether it produces an inspectable policy. All
numbers are reported as normalized liability offsets, not complete economic
P\&L. Consequently, a value closer to zero is better for the mean and lower-tail
statistics, but zero is not a break-even threshold. Comparisons are made under
the same liability, hedge instruments, costs, and bounds. The numerical tables
report the current panel estimates and are interpreted descriptively; the
protocol-locked rerun will add the paired circular-block intervals specified in
Section~\ref{subsec:evaluation_metrics} before inferential claims are made.

\subsection{European-Option Panel Results}
\label{subsec:european_results}

Table~\ref{tab:european-panel-average} gives the equal-ticker average for the
European short-call task. The two ProtoHedge rows correspond to configurations
chosen by the prespecified mean- and tail-oriented validation rules. This is
important: they represent two ex ante decision objectives, not two models
chosen after viewing the test metric.

\begin{table*}[t]
\centering
\small
\begin{tabular}{lrrr}
\toprule
Model & Mean offset & Offset CVaR$_{5\%}$ & Bound occupancy \\
\midrule
Unhedged & -0.030502 & -0.140291 & 0.000000 \\
Spot-Delta & -0.030331 & -0.137235 & 0.000000 \\
Spot-Delta Band & -0.030295 & -0.137247 & 0.000000 \\
Deep Hedging & -0.030277 & -0.134793 & 0.077515 \\
ProtoHedge (Validation-Mean) & \textbf{-0.029651} & -0.133763 & 0.048616 \\
ProtoHedge (Validation-Tail) & -0.030145 & \textbf{-0.132121} & \textbf{0.030434} \\
\bottomrule
\end{tabular}
\caption{European short-call results, averaged across seeds within ticker and
then equally across ten tickers. Higher mean offset and CVaR$_{5\%}$ are
better. Bound occupancy is the fraction of path--time--instrument inventories
at a hard cumulative-position limit; lower is better among learned policies.}
\label{tab:european-panel-average}
\end{table*}

The first result is that the neural policies improve lower-tail offset relative
to the transparent baselines. Deep Hedging increases average CVaR$_{5\%}$ from
$-0.137235$ under Spot-Delta to $-0.134793$. ProtoHedge improves it further:
the Validation-Mean policy reaches $-0.133763$, while the Validation-Tail policy
reaches $-0.132121$. Thus the interpretation constraint does not produce a
visible panel-level tail penalty in this experiment.

The mean-oriented ProtoHedge policy also has the highest average offset. Its
improvement over Deep Hedging is $0.000625$, or approximately $2.07\%$ relative
to the magnitude of the Deep Hedging mean offset. The tail-oriented policy
sacrifices some of that mean advantage but remains $0.000132$ above Deep
Hedging while improving CVaR$_{5\%}$ by $0.002672$, approximately $1.98\%$ of
the Deep Hedging tail-loss magnitude. The two selections therefore expose a
small but economically coherent mean--tail tradeoff rather than a universal
winner.

Table~\ref{tab:european-ticker} checks whether these averages are driven by a
small subset of assets. The Validation-Mean policy has a higher mean offset
than Deep Hedging for all ten tickers. The Validation-Tail policy has a higher
CVaR$_{5\%}$ for all ten. The absolute improvements are largest for the most
volatile single-name series, NVDA, but the direction also holds for broad
market, sector, and Treasury ETFs.

\begin{table*}[t]
\centering
\small
\setlength{\tabcolsep}{4.5pt}
\begin{tabular}{lrrrrrr}
\toprule
& \multicolumn{2}{c}{Mean offset} & \multicolumn{2}{c}{CVaR$_{5\%}$} &
\multicolumn{2}{c}{Bound occupancy} \\
\cmidrule(lr){2-3}\cmidrule(lr){4-5}\cmidrule(lr){6-7}
Ticker & DH & Proto-Mean & DH & Proto-Tail & DH & Proto-Tail \\
\midrule
AAPL & -0.0396 & -0.0379 & -0.1479 & -0.1464 & 0.1250 & 0.0096 \\
IWM  & -0.0271 & -0.0270 & -0.1240 & -0.1225 & 0.0907 & 0.0028 \\
NVDA & -0.0724 & -0.0699 & -0.3265 & -0.3204 & 0.1744 & 0.0116 \\
QQQ  & -0.0288 & -0.0286 & -0.1104 & -0.1069 & 0.1660 & 0.0304 \\
SPY  & -0.0217 & -0.0216 & -0.0859 & -0.0840 & 0.1927 & 0.0001 \\
TLT  & -0.0143 & -0.0134 & -0.0819 & -0.0811 & 0.0000 & 0.0835 \\
XLE  & -0.0286 & -0.0285 & -0.1705 & -0.1688 & 0.0000 & 0.0151 \\
XLF  & -0.0238 & -0.0237 & -0.1048 & -0.1021 & 0.0000 & 0.1361 \\
XLK  & -0.0281 & -0.0277 & -0.1093 & -0.1064 & 0.0264 & 0.0119 \\
XLV  & -0.0184 & -0.0182 & -0.0866 & -0.0827 & 0.0000 & 0.0032 \\
\bottomrule
\end{tabular}
\caption{Ticker-level European comparison. Proto-Mean is the
validation-mean selection and Proto-Tail is the validation-tail selection.
The paired columns deliberately show the criterion targeted by each selector.
Higher offsets and lower occupancy are preferable.}
\label{tab:european-ticker}
\end{table*}

These differences are modest, which is consistent with the objective of the
paper. ProtoHedge is not intended to obtain large gains by solving a different
hedging problem; it is intended to retain the quality of a black-box policy
while exposing a finite set of representative states. The relevant empirical
finding is therefore that the prototype bottleneck remains competitive across
the panel, not that it creates a new source of arbitrage.

\subsection{Position-Bound Analysis}
\label{subsec:bound_results}

The performance result is more credible if it is not generated by persistent
saturation at the hard inventory limits. In the European task, average bound
occupancy is $0.0775$ for Deep Hedging, $0.0486$ for Validation-Mean
ProtoHedge, and $0.0304$ for Validation-Tail ProtoHedge. These correspond to
average reductions of approximately $37.3\%$ and $60.7\%$, respectively. The
prototype policies therefore achieve their panel-level mean and tail results
while operating more often in the interior of the feasible set.

The ticker-level pattern is not uniform. The tail-oriented ProtoHedge policy
has lower occupancy than Deep Hedging for six of ten underlyings, including
large reductions for AAPL, IWM, NVDA, QQQ, and SPY. It has higher occupancy for
TLT, XLE, XLF, and XLV, where the Deep Hedging benchmark is at zero or close to
it. This heterogeneity is why occupancy is reported beside performance rather
than used as a stand-alone measure of model quality. A lower-bound policy is
not automatically a better hedge; it is favorable only when offset quality is
preserved.

The Asian task produces a different boundary pattern, discussed below. This
cross-liability contrast is informative: the prototype bottleneck appears to
regularize inventory in the European experiment, but path dependence can
require a different use of the feasible action region.

\subsection{Interpretability Verification: QQQ Case Study}
\label{subsec:interpretability_results}

Interpretability requires evidence about the prototype mechanism itself, not
only competitive performance. For a detailed case study, we examine policies
for the Invesco QQQ Trust (QQQ), a highly traded exchange-traded fund that
tracks the Nasdaq-100. QQQ provides a long and liquid option history and has
complete representative artifacts for both liabilities, allowing the effect
of path dependence on prototype use to be studied on the same underlying. It
is used for interpretability illustration, not because it had the strongest
test performance; panel-level performance claims continue to use all available
tickers. Table~\ref{tab:prototype-concentration}
reports cumulative prototype-assignment mass and two effective-library-size
statistics. The entropy effective size is
\[
K_{\mathrm{ent}}=\exp\!\left(-\sum_{k=1}^{K}\bar\alpha_k
\log\bar\alpha_k\right),
\]
and the inverse-concentration size is
\[
K_{\mathrm{inv}}=\left(\sum_{k=1}^{K}\bar\alpha_k^2\right)^{-1},
\]
where $\bar\alpha_k$ is the prototype's mean assignment weight over test paths
and decision times. Both equal $K$ under perfectly uniform use and approach
one as assignment mass concentrates.

\begin{table*}[t]
\centering
\small
\begin{tabular}{llrrrrrrr}
\toprule
Liability & Representative policy & $K$ & Top 1 & Top 3 & Top 5 & Top 10 &
$K_{\mathrm{ent}}$ & $K_{\mathrm{inv}}$ \\
\midrule
European & QQQ ProtoHedge & 25 & 0.208 & 0.584 & 0.747 & 0.996 & 8.29 & 7.04 \\
Asian & QQQ ProtoHedge & 50 & 0.199 & 0.352 & 0.480 & 0.713 & 18.78 & 13.04 \\
\bottomrule
\end{tabular}
\caption{Concentration of prototype assignment mass for the representative QQQ
case study. Top-$k$ columns report cumulative mass after sorting prototypes by
mean use. Effective sizes below the nominal $K$ indicate non-uniform use.}
\label{tab:prototype-concentration}
\end{table*}

The European model uses a compact effective dictionary. Three of its 25
prototypes explain $58.4\%$ of assignment mass, ten explain $99.6\%$, and its
effective size is between seven and eight. The prototype layer is therefore
not a decorative uniform mixture: most decisions pass through a substantially
smaller set of recurring states. The Asian model uses more of its available
library. Its top ten prototypes explain $71.3\%$ and its effective size is
between 13 and 19. This is still concentrated relative to 50 uniform
prototypes, but it suggests that a path-dependent payoff requires a richer set
of state representatives.

Table~\ref{tab:prototype-regimes} reports the most heavily weighted prototype
within return, volatility, and drawdown terciles. In the European model,
prototype 21 dominates down-return, high-volatility, and mild-drawdown states,
whereas prototype 8 dominates middle/up returns, low/middle volatility, and
middle/deep drawdowns. This systematic switch is consistent with regime-aware
example use. In the Asian model, prototype 10 remains the single most-used
anchor in every tercile, but its weight falls from $0.215$ in low volatility to
$0.182$ in high volatility while the remaining assignment mass redistributes.
The Asian evidence therefore supports state-dependent weighting, but not the
stronger claim that the identity of the top prototype changes in every regime.

\begin{table*}[t]
\centering
\small
\begin{tabular}{llccc}
\toprule
Liability & Regime partition & Low / Down / Mild & Middle & High / Up / Deep \\
\midrule
European & Return & P21 (0.206) & P8 (0.250) & P8 (0.242) \\
European & Volatility & P8 (0.313) & P8 (0.203) & P21 (0.208) \\
European & Drawdown & P21 (0.210) & P8 (0.225) & P8 (0.295) \\
Asian & Return & P10 (0.199) & P10 (0.204) & P10 (0.193) \\
Asian & Volatility & P10 (0.215) & P10 (0.200) & P10 (0.182) \\
Asian & Drawdown & P10 (0.190) & P10 (0.200) & P10 (0.206) \\
\bottomrule
\end{tabular}
\caption{Dominant prototype and mean assignment weight in the QQQ case study.
Regimes are test-period terciles defined separately for terminal return,
average implied volatility, and maximum drawdown.}
\label{tab:prototype-regimes}
\end{table*}

Nearest-state retrieval supplies the final link from model internals to market
states. Table~\ref{tab:nearest-states} reports the closest observed QQQ state
for four highly used European prototypes. For example, P8 represents a state
with normalized spot above inception and a moderate positive spot inventory,
whereas P21 is associated with normalized spot near $0.85$, a more expensive
listed call, and a larger spot inventory. P11 and P16 occupy intermediate
price and inventory regions. These are concrete observed states, not latent
centroids. The retrieved path identifiers allow the complete 20-step histories
to be inspected and, in the locked rerun, mapped to calendar dates.

\begin{table*}[t]
\centering
\small
\begin{tabular}{lrrrrrr}
\toprule
Prototype & Path & Step & Spot position & Option position & Norm. spot & Norm. option \\
\midrule
P8  & 315 & 11 & 0.5612 & -0.0012 & 1.0464 & 0.0183 \\
P11 & 497 & 5  & 0.4701 & -0.0013 & 0.9940 & 0.0600 \\
P16 & 36  & 5  & 0.5232 & -0.0011 & 1.0129 & 0.0255 \\
P21 & 372 & 7  & 0.6582 & -0.0018 & 0.8492 & 0.0673 \\
\bottomrule
\end{tabular}
\caption{Nearest observed states in the QQQ case study for four dominant
European prototypes. Prices are normalized by episode-initial spot. These
state summaries support example-based inspection; calendar-event labels will
be added from the contract-consistent rerun rather than inferred from path
shape alone.}
\label{tab:nearest-states}
\end{table*}

Together, concentration, regime use, and nearest-state retrieval verify three
different aspects of interpretability: the policy uses a finite and non-uniform
state dictionary, the assignment mechanism responds to market conditions, and
dominant anchors can be inspected as observed states. They do not yet establish
that professional hedgers find the explanations operationally useful, nor do
they justify labeling a retrieved state as a crisis or flash crash without its
calendar metadata.

\subsection{Asian-Option Results}
\label{subsec:asian_results}

The Asian experiment changes the liability and adds running-average moneyness
to every learned policy and prototype state, while preserving the data split,
hedge instruments, costs, candidate grid, and selection rules. It therefore
tests whether ProtoHedge remains competitive when current prices alone are not
a sufficient liability state.

\begin{table*}[t]
\centering
\small
\begin{tabular}{lrrrr}
\toprule
Model & Mean offset & Offset CVaR$_{5\%}$ & Bound occupancy & Tickers \\
\midrule
Unhedged & -0.016479 & -0.078681 & 0.000000 & 10 \\
Spot-Delta & -0.016265 & -0.074999 & 0.000000 & 10 \\
Spot-Delta Band & -0.016240 & -0.074893 & 0.000000 & 10 \\
Deep Hedging & -0.015801 & -0.074384 & 0.018519 & 9 \\
ProtoHedge (Validation-Mean) & \textbf{-0.015530} & -0.073844 & 0.019606 & 9 \\
ProtoHedge (Validation-Tail) & -0.015786 & \textbf{-0.071501} & 0.038187 & 9 \\
\bottomrule
\end{tabular}
\caption{Asian short-call results averaged across the available ticker runs.
The final column reports the number of tickers contributing to each row.}
\label{tab:asian-panel-average}
\end{table*}

The principal performance conclusion survives path dependence. Relative to
Deep Hedging, the mean-selected policy improves average mean offset by
$0.000271$ and CVaR$_{5\%}$ by $0.000541$. It wins seven of nine available
ticker comparisons on each metric. The tail-selected policy has essentially
the same mean offset as Deep Hedging, a difference of $0.000015$, while
improving average CVaR$_{5\%}$ by $0.002883$, or approximately $3.88\%$ of the
Deep Hedging tail-loss magnitude. Its CVaR is higher for all nine available
tickers. This is the strongest current evidence that the prototype bottleneck
remains useful when the liability is path-dependent.

The position result does not carry over unchanged. The mean-selected policy
has average occupancy $0.0196$ versus $0.0185$ for Deep Hedging, and the
tail-selected policy reaches $0.0382$. Unlike the European task, neither Asian
selection lowers occupancy ticker by ticker. The path-dependent liability can
therefore preserve the performance--interpretability tradeoff without
preserving the same inventory-regularization pattern. This distinction is more
informative than a blanket claim that ProtoHedge always trades less
aggressively.

\begin{table*}[t]
\centering
\small
\setlength{\tabcolsep}{4.5pt}
\begin{tabular}{lrrrrrr}
\toprule
& \multicolumn{2}{c}{Mean offset} & \multicolumn{2}{c}{CVaR$_{5\%}$} &
\multicolumn{2}{c}{Bound occupancy} \\
\cmidrule(lr){2-3}\cmidrule(lr){4-5}\cmidrule(lr){6-7}
Ticker & DH & Proto-Mean & DH & Proto-Tail & DH & Proto-Tail \\
\midrule
IWM  & -0.0154 & -0.0151 & -0.0751 & -0.0715 & 0.0000 & 0.0000 \\
NVDA & -0.0375 & -0.0364 & -0.1672 & -0.1635 & 0.1667 & 0.1949 \\
QQQ  & -0.0154 & -0.0156 & -0.0627 & -0.0583 & 0.0000 & 0.0000 \\
SPY  & -0.0120 & -0.0120 & -0.0532 & -0.0519 & 0.0000 & 0.0573 \\
TLT  & -0.0080 & -0.0073 & -0.0449 & -0.0427 & 0.0000 & 0.0068 \\
XLE  & -0.0154 & -0.0153 & -0.0882 & -0.0855 & 0.0000 & 0.0814 \\
XLF  & -0.0133 & -0.0132 & -0.0643 & -0.0616 & 0.0000 & 0.0014 \\
XLK  & -0.0149 & -0.0148 & -0.0618 & -0.0591 & 0.0000 & 0.0007 \\
XLV  & -0.0103 & -0.0101 & -0.0522 & -0.0493 & 0.0000 & 0.0012 \\
\bottomrule
\end{tabular}
\caption{Ticker-level Asian comparison across the available learned-policy
runs.}
\label{tab:asian-ticker}
\end{table*}

Table~\ref{tab:cross-liability-summary} summarizes the central comparison with
Deep Hedging. Across both liabilities, the validation-tail policy provides the
most consistent lower-tail improvement. The mean selector provides the larger
average-offset gain. The European and Asian occupancy columns then show why
performance and constraint use must be assessed jointly rather than collapsed
into one score.

\begin{table*}[t]
\centering
\small
\begin{tabular}{llrrrrrr}
\toprule
Liability & Selection & $N$ & $\Delta$Mean & Mean wins & $\Delta$CVaR & CVaR wins & $\Delta$Bound \\
\midrule
European & Validation-Mean & 10 & 0.000625 & 10 & 0.001030 & 7 & -0.028898 \\
European & Validation-Tail & 10 & 0.000132 & 6 & 0.002672 & 10 & -0.047081 \\
Asian & Validation-Mean & 9 & 0.000271 & 7 & 0.000541 & 7 & 0.001087 \\
Asian & Validation-Tail & 9 & 0.000015 & 5 & 0.002883 & 9 & 0.019669 \\
\bottomrule
\end{tabular}
\caption{ProtoHedge minus Deep Hedging, averaged equally across available
tickers. Positive mean and CVaR differences favor ProtoHedge; negative bound
differences favor ProtoHedge. ``Wins'' counts tickers with a favorable paired
point estimate. Statistical claims require the locked block-bootstrap
intervals.}
\label{tab:cross-liability-summary}
\end{table*}

Overall, the Asian experiment extends rather than merely repeats the European
finding. The learned policy receives the payoff state that makes the Asian
problem Markov in the supplied information, the tail-oriented prototype policy
remains competitive on every available ticker, and its prototype library is
inspectable but less concentrated. At the same time, higher boundary use shows
that path dependence changes the economic behavior of the selected policy.
\section{Discussion and Limitations}
\label{sec:discussion}

The results answer the paper's central empirical question in qualified but
favorable terms. A fixed prototype bottleneck can remain competitive with a
common-protocol Deep Hedging network on historical data. In the European task,
the mean-oriented and tail-oriented selectors improve the corresponding panel
averages and reduce average use of hard position limits. In the Asian task,
the same selectors preserve mean and lower-tail competitiveness after adding
the running payoff state, but the reduction in boundary use disappears. The
interpretability evidence also changes with the liability: the European policy
uses a compact effective dictionary and switches its dominant anchor across
several regimes, whereas the Asian policy distributes mass across a larger
library while retaining one dominant prototype.

\subsection{Economic and Modeling Interpretation}

The magnitude of the performance differences should be interpreted in the
context of the research objective. All learned models trade the same two
instruments, observe the same task-relevant information, optimize the same
monetary utility, and face the same costs and bounds. ProtoHedge's contribution
is therefore not a fundamentally different source of hedging returns. It is a
restriction on how a neural action is represented. Preserving or slightly
improving offset quality under that restriction is evidence that the black-box
benchmark may use more representational freedom than this empirical task
requires.

The separate validation selectors are also informative. The mean-oriented
configuration yields the larger average-offset improvement, while the
tail-oriented configuration yields the more consistent CVaR improvement. This
means that ``best ProtoHedge'' is not a single economically neutral concept.
The appropriate model depends on whether a decision maker prioritizes average
liability offset or severe-loss mitigation. The prototype layer makes this
choice easier to audit because the selected policy can be examined through its
prototype actions, assignment concentration, and representative states.

Position-bound occupancy adds a behavioral dimension that is invisible in
mean and CVaR alone. Lower European occupancy suggests that the prototype
restriction can act as an inventory regularizer. The Asian result prevents us
from generalizing that claim: path-dependent hedging uses the feasible region
more aggressively even when performance remains competitive. The practical
lesson is that interpretability, tail protection, and conservative inventory
use are related but distinct properties. They should be reported jointly, not
combined into a single ranking.

\subsection{Limitations and Cases Without Dominance}

ProtoHedge does not dominate every model, asset, or criterion. The
mean-oriented and tail-oriented policies are different validation selections;
neither is claimed to optimize all metrics simultaneously. European occupancy
is higher than Deep Hedging for several individual tickers even though the
panel average is lower, and both Asian selections have higher average
occupancy. Some per-ticker mean comparisons also favor Deep Hedging. These
exceptions are part of the empirical result rather than failures to be hidden:
they identify settings in which the prototype restriction or the selected
objective trades flexibility for auditability.

Several broader limitations remain. First, the reported outcome excludes the
initial liability premium. It is a liability-offset measure for controlled
policy comparison, not a complete terminal P\&L, and zero does not represent
economic break-even. A complete P\&L study would require an observed or
independently priced liability premium under a clearly stated valuation rule.
This is also why the negative-offset frequency is not a headline metric.

Second, the transaction-cost formula is explicit and applied consistently but
is stylized. The backtest does not reconstruct executable bid and ask trades,
market impact, funding, margin, assignment risk, or desk-level risk limits. It
also assumes that the selected listed hedge option remains tradable at the
observed daily midpoint throughout an episode. Contract continuity corrects
the self-financing identity, but it does not eliminate execution uncertainty.

Third, the panel is broad relative to a single-ticker case study but remains
limited to liquid U.S. equities and ETFs. It does not establish external
validity for index options, foreign exchange, commodities, rates, digital
assets, or less liquid single names. The 20-step horizon and near-ATM call
selection similarly cover only one part of the option surface.

Fourth, windows overlap within each calendar period. Splitting before window
construction and imposing a 19-observation embargo eliminate shared dates
across train, validation, and test sets, but they do not make adjacent test
windows independent. The final claims must therefore use the paired
circular-block intervals specified in Section~\ref{subsec:evaluation_metrics},
not standard errors that treat all episodes as independent. The current point
estimates should not be described as statistically significant until those
intervals are populated.

Fifth, the final protocol-locked run is required to regenerate every table
under cumulative
inventory accounting, chronological purged splits, fixed hedge contracts,
payoff-sufficient Asian state, and validation-only configuration selection.
The paper is structured so those results replace numerical cells without
changing definitions or claims logic.

Sixth, inspectability is not the same as practitioner validation. Concentrated
weights and nearest historical states show what the model uses, but the present
study does not ask professional hedgers whether those explanations improve
model review or trading decisions. Likewise, the retrieved-state artifact does
not yet attach calendar dates and event annotations. We therefore do not claim
that the current prototypes have already been verified as financial-crisis or
flash-crash states. Date-aware retrieval and a prespecified event study are
needed for that stronger claim.

Finally, the benchmark set is intentionally focused: transparent delta rules
and a common-protocol feed-forward Deep Hedging network. It does not include the
full range of recurrent, transformer, distributional, or dynamics-aware deep
hedging methods, including Deeper Hedging under Chiarella--Heston dynamics
\cite{gao2023deeper}. Nor does the paper complete the original scalability
agenda. The grid varies $K$, source policy, and similarity weighting, but it
does not systematically isolate wall-clock complexity, memory use, feature
dimension, or approximate nearest-neighbor search. Those are distinct
comparative and systems studies rather than evidence already supplied by this
run.

\subsection{Relation to the ProtoHedge Research Agenda}

The original ProtoHedge paper identified empirical validation,
path-dependent instruments, historical-scenario interpretation, scalability,
and broader model comparison as future work \cite{faloughi2025protohedge}. The
present study directly advances the first two items: it constructs a real-data
panel, enforces historical out-of-sample evaluation, and adds an Asian
liability with a payoff-sufficient state. It partially advances historical
interpretation by measuring concentration, regime use, and nearest-state
retrieval. It does not claim to complete the professional-user, scalability,
or broad-benchmark components.

This distinction keeps the paper's contribution clear. A single additional
run cannot credibly ``solve'' every future-work axis because each requires a
different experimental design. Scalability requires controlled timing and
memory experiments as dimensionality and $K$ vary. Comparative evaluation
requires faithful implementation and tuning of additional models. Practitioner
validation requires a human-subject or structured expert protocol. These are
natural follow-on projects, but they should not dilute the central contribution
of the present empirical test.

\section{Conclusion}

This paper tests whether ProtoHedge's fixed, example-based representation
survives the transition from simulation to historical option data. The
framework uses one consistent listed hedge contract per episode, chronological
purged splits, and cumulative-inventory accounting. It also imposes stylized
costs, hard action and position bounds, validation-only selection, and a
payoff-sufficient state for the Asian liability. The reported outcome excludes
the liability premium and therefore is not complete economic P\&L.

Across the current panel estimates, ProtoHedge remains competitive with Deep
Hedging on both liabilities. The European mean and tail selections improve
their targeted average offset metrics while reducing average bound occupancy.
The Asian selections preserve the performance result, particularly in the
lower tail, but use the hard bounds more frequently. Prototype assignments are
non-uniform, respond to market regimes, and map back to observed historical
states; the Asian policy uses a larger effective dictionary than the European
policy. These findings support the original ProtoHedge thesis in a qualified
empirical form: interpretability can be introduced at little apparent cost to
hedging quality, but its effects on inventory behavior and prototype
concentration depend on the liability.

The remaining work before submission is finite and explicit: complete the
protocol-locked rerun, populate paired circular-block confidence intervals,
attach dates to nearest historical states, and update the numerical cells
without changing the prespecified analysis. Subject to those checks, the paper
provides the first
systematic real-data validation of the fixed-prototype hedging framework and a
foundation for subsequent work on scalable comparison and trainable prototype
locations.
\section*{Ethics and Privacy Statement}

This study uses licensed historical market data and does not involve human
subjects or personal data. A principal risk is that backtest results could be
misinterpreted as evidence of live-trading profitability; we therefore state
the premium convention, execution assumptions, uncertainty, and external
validity limitations explicitly. Prototype explanations may improve model
auditability, but they should support rather than replace professional risk
management, independent validation, and applicable regulatory controls.

%%
%% The next two lines define the bibliography style to be used, and
%% the bibliography file.
%%
%% The next two lines define the bibliography style to be used, and
%% the bibliography file.
\bibliographystyle{ACM-Reference-Format}
\bibliography{references}

\appendix

\section{Outcome and Inference Details}
\label{app:inference}

For a set of test offsets $\{\mathcal{O}_i\}_{i=1}^{n}$, let
$q_{0.05}$ be the empirical fifth percentile. The reported lower-tail
statistic is
\[
\widehat{\mathrm{CVaR}}_{5\%}
=
\frac{1}{|\mathcal{I}_{0.05}|}
\sum_{i\in\mathcal{I}_{0.05}}\mathcal{O}_i,
\qquad
\mathcal{I}_{0.05}=\{i:\mathcal{O}_i\leq q_{0.05}\}.
\]
Because the outcome is a gain-like offset, a larger value is preferable. For
RMSE and downside deviation, the comparison signs are reversed before
reporting a ProtoHedge improvement so that every paired effect has the same
orientation.

For each ticker, liability, selector, and test episode, offsets are first
averaged across available training seeds. ProtoHedge and Deep Hedging are then
kept paired on the same episode indices. A circular block sample concatenates
random blocks of 20 adjacent episode indices, wrapping at the end of the test
sequence, until the original sample length is reached. The paired statistic is
recomputed for each of 2,000 replicates. Panel replicates resample within every
ticker independently and then take the equal-ticker average. Percentile 2.5\%
and 97.5\% quantiles define the reported interval.

\section{Data-Integrity Checks}
\label{app:data_checks}

The final data artifact records the following checks for every retained
sample:
\begin{enumerate}
    \item all 20 observations belong to one chronological split and one
    contiguous date block;
    \item the 19-observation embargo is applied after each split boundary;
    \item one \texttt{optionid}, strike, and expiration are used throughout
    the episode;
    \item the listed hedge option remains unexpired after the liability
    horizon and passes the initial liquidity and moneyness filters;
    \item one-step listed-option returns are multiplied by cumulative option
    inventory, not only the newest trade;
    \item Asian policies and prototype features contain running average
    moneyness at every decision time; and
    \item model configuration files identify whether an artifact was produced
    by the corrected scientific protocol, preventing reuse of superseded
    random-split or contract-switching results.
\end{enumerate}

\section{Recommended Supplementary Reporting}
\label{app:supplementary}

The main paper emphasizes compact tables. The accompanying supplementary
artifact should archive, without adding visual clutter to the core argument:
\begin{enumerate}
    \item full per-ticker results for every classical baseline, Deep Hedging,
    and both validation-selected ProtoHedge policies;
    \item validation-only candidate tables over prototype count, source
    policy, and similarity weighting;
    \item per-ticker paired circular-block intervals for mean offset,
    CVaR$_{5\%}$, RMSE, and downside deviation;
    \item prototype actions, assignment frequencies, concentration curves,
    and regime-conditioned assignment tables for each selected model; and
    \item nearest historical states with ticker, date, contract identifier,
    strike, expiration, payoff state, and complete 20-step price trajectory.
\end{enumerate}

These materials separate model selection from final test reporting and permit
readers to inspect the evidence underlying every aggregate claim.

\end{document}
\endinput
%%
%% End of file `Final_Paper.tex'.
